% ============================================================
%  VI. CONCLUSION
% ============================================================
\section{Conclusion}
\label{sec:conclusion}

This paper presented \ac{TERESA}, a legged mobile manipulator for autonomous
emergency assessment. The core contribution is a whole-body active perception
framework that treats body posture---height, pitch, and yaw---as a perceptual
resource, closing the perception-action loop at the whole-body level.
Experiments on the physical Spot+Z1 platform demonstrated that the arm's
\SI{360}{\degree} coverage accelerates search by a factor of five over a
fixed-camera baseline, the dual-sensor lock prevents single-camera detection
failures, anticipatory scanning eliminates idle time by completing the body
scan concurrently with approach, and body reconfiguration lifts anatomical
target completion from \SI{40}{\percent} to \SI{100}{\percent} for \ac{FAST}
targets.  The exposure pipeline integrates zero-shot wound detection with an
interactive web-based review interface.

Several limitations remain: the system assumes a single stationary mannequin on
flat terrain; ultrasound image quality and clinical scanning pressure were not
evaluated; and injury detection relies on pre-trained models whose performance
may vary across wound types and skin tones.  Future work will extend
\ac{TERESA} to multi-patient scenarios, clinical ultrasound acquisition with
impedance-controlled probe contact, and ultimately the full \ac{ABCDE}
protocol, providing a complete autonomous primary survey for emergency triage.
